\chapter{Persepuluhan dan Perseratusan}\label{ch:g4:decimals}

Di antara $3$ dan $4$ bilangan bulat membisu --- padahal seorang
pelari bisa unggul $3$ meter lebih sedikit. Memotong satuan menjadi
sepuluh, lalu seratus, menghasilkan \emph{\omterm{def:g4:fractions:def}{pecahan} desimal}, beserta
singkatan yang menakjubkan untuknya: \omterm{def:g4:decimals:point}{titik desimal}. (Bilangan
desimal mendapat babnya sendiri pada \cref{ch:g5:decimals}.)

\section{Pecahan desimal}

\begin{definition}[Persepuluhan dan perseratusan]\label{def:g4:decimals:def}
Memotong satu satuan menjadi $10$ bagian yang sama menghasilkan
\emph{persepuluhan}, $\frac{1}{10}$; memotong setiap persepuluhan
menjadi $10$ lagi menghasilkan \emph{perseratusan}, $\frac{1}{100}$.
Jadi
\[
1 = \frac{10}{10} = \frac{100}{100},
\qquad
\frac{1}{10} = \frac{10}{100}.
\]
\end{definition}

\begin{omfigure}
\begin{tikzpicture}[scale=0.42]
  \draw[gray!50, thin] (0,0) grid (10,10);
  \foreach \i in {0,...,9} \fill[omDef!35] (3,\i+0.06)
    rectangle (3.94,\i+0.94);
  \fill[omThm!70] (6.05,4.05) rectangle (6.95,4.95);
  \draw[thick] (0,0) rectangle (10,10);
  \node[below, font=\small] at (5,-0.4)
    {seluruh persegi $= 1$;\ \ satu kolom (biru) $= \frac{1}{10}$;\ \
    satu petak kecil (merah) $= \frac{1}{100}$};
\end{tikzpicture}

{\small Persegi seratus: sepuluh kolom berisi sepuluh petak. Satu
kolom adalah satu persepuluhan, satu petak satu perseratusan.}
\end{omfigure}

\begin{example}\label{ex:g4:decimals:count}
Mewarnai $3$ kolom dan $7$ petak tambahan pada persegi seratus berarti mewarnai
\[
\frac{3}{10} + \frac{7}{100} = \frac{30}{100} + \frac{7}{100}
= \frac{37}{100}
\]
bagian persegi itu --- tiga puluh tujuh perseratusan.
\end{example}

\section{Titik desimal}

\begin{definition}[Penulisan desimal]\label{def:g4:decimals:point}
Bilangan $2 + \frac{3}{10} + \frac{7}{100}$ ditulis
\[
2.37
\]
--- \emph{titik desimal}\index{titik desimal} memisahkan bagian bulat ($2$) dari
bagian desimalnya; angka pertama sesudah titik membilang
persepuluhan, angka kedua membilang perseratusan. Dibaca: ``dua
titik tiga tujuh'', atau ``dua dan tiga puluh tujuh perseratus''.
\end{definition}

\begin{example}\label{ex:g4:decimals:convert}
\[
\frac{7}{10} = 0.7, \qquad
\frac{43}{100} = 0.43, \qquad
\frac{5}{100} = 0.05, \qquad
3 + \frac{4}{10} = 3.4 .
\]
Hati-hati dengan $\frac{5}{100}$: lima \emph{perseratusan}
memerlukan angka $0$ di tempat persepuluhan --- $0.05$, bukan $0.5$.
\end{example}

\begin{omfigure}
\begin{tikzpicture}[scale=1.35]
  \draw[-Stealth] (-0.2,0) -- (8.4,0);
  \foreach \i in {0,10,20,30,40,50,60,70,80}
    \draw ({\i*0.1},-0.09) -- ({\i*0.1},0.09);
  \foreach \i in {0,...,80} \draw ({\i*0.1},-0.045) -- ({\i*0.1},0.045);
  \foreach \x/\l in {0/2, 1/{2.1}, 2/{2.2}, 3/{2.3}, 4/{2.4},
                     5/{2.5}, 6/{2.6}, 7/{2.7}, 8/{2.8}}
    \node[below, font=\footnotesize] at (\x,-0.1) {$\l$};
  \fill[omThm] (3.7,0) circle (1.7pt)
    node[above, font=\small] {$2.37$};
\end{tikzpicture}

{\small Memperbesar tampilan antara $2$ dan $2.8$: tanda besar
setiap persepuluhan, tanda kecil setiap perseratusan. Bilangan
$2.37$ terletak antara $2.3$ dan $2.4$, pada tanda kecil ke-$7$.}
\end{omfigure}

\begin{example}[Uang adalah perseratusan]\label{ex:g4:decimals:money}
Satu sen adalah seperseratus satu euro (atau dolar, atau pound):
harga $4.85$ berarti $4$ satuan penuh dan $85$ perseratusan. Uang
adalah \omterm{def:g4:decimals:point}{titik desimal} yang setiap hari sudah kamu pakai: $10$ sen
$= 0.10$; dua euro lima puluh $= 2.50$; dan $0.05$ adalah uang logam
lima sen, bukan lima puluh!
\end{example}

\section{Membandingkan dengan persepuluhan}

\begin{method}[Membandingkan bilangan desimal (perkenalan)]\label{met:g4:decimals:compare}
\begin{enumerate}
  \item Bandingkan bagian bulatnya dulu: $3.2 > 2.9$.
  \item Bagian bulatnya sama: bandingkan persepuluhannya, lalu
        perseratusannya: $2.37 < 2.5$ karena $3$ persepuluhan
        $< 5$ persepuluhan.
  \item Menulis keduanya dengan dua angka desimal itu membantu:
        $2.5 = 2.50$, dan $37 < 50$ perseratusan.
\end{enumerate}
Perangkapnya: $2.37$ punya angka lebih banyak daripada $2.5$, tetapi
\emph{lebih kecil} --- angka yang lebih banyak tidak berarti lebih besar!
\end{method}

\begin{example}\label{ex:g4:decimals:compare}
Urutkan $0.4$, $0.35$ dan $0.09$: dengan perseratusan, $0.40$,
$0.35$, $0.09$; jadi
\[
0.09 < 0.35 < 0.4 .
\]
\end{example}

\section{Latihan}

\begin{exercise}[$\star$]\label{exo:g4:decimals:1}
Pada persegi seratus, berapa bagian yang terwarnai jika kita
mewarnai: $5$ kolom; $2$ kolom dan $3$ petak; $45$ petak? Jawablah dengan \omterm{def:g4:fractions:def}{pecahan}.
\end{exercise}

\begin{exercise}[$\star$]\label{exo:g4:decimals:2}
Tulislah dengan \omterm{def:g4:decimals:point}{titik desimal}:
$\frac{9}{10}$; \quad $\frac{27}{100}$; \quad $\frac{3}{100}$; \quad
$5 + \frac{6}{10}$; \quad $\frac{60}{100}$.
\end{exercise}

\begin{exercise}[$\star$]\label{exo:g4:decimals:3}
Tulislah sebagai \omterm{def:g4:fractions:def}{pecahan} (persepuluhan atau perseratusan):
$0.3$; \quad $0.51$; \quad $0.07$; \quad $1.9$.
\end{exercise}

\begin{exercise}[$\star$]\label{exo:g4:decimals:4}
Pada $6.58$: apa yang dibilang angka $5$? Angka $8$? Uraikan
bilangan itu seperti pada \cref{def:g4:decimals:point}.
\end{exercise}

\begin{exercise}[$\star$]\label{exo:g4:decimals:5}
Letakkan pada garis bilangan berskala persepuluhan dari $4$ sampai
$5$: $4.2$; \ $4.75$ (di antara dua tanda yang mana?); \ $4.9$.
\end{exercise}

\begin{exercise}[$\star$]\label{exo:g4:decimals:6}
Salin dan lengkapi dengan $<$, $>$ atau $=$:
\[
0.6 \;?\; 0.60, \qquad
2.8 \;?\; 2.75, \qquad
0.09 \;?\; 0.1, \qquad
5.3 \;?\; 5.13 .
\]
\end{exercise}

\begin{exercise}[$\star$]\label{exo:g4:decimals:7}
Urutkan dari yang terkecil ke yang terbesar: $1.05$; \ $1.5$; \ $0.95$; \ $1.45$.
\end{exercise}

\begin{exercise}[$\star$]\label{exo:g4:decimals:8}
Tulislah dalam penulisan desimal: tiga euro lima sen; dua belas euro
lima puluh sen; sembilan puluh sembilan sen. Lalu urutkan ketiga harga itu.
\end{exercise}

\begin{exercise}[$\star$]\label{exo:g4:decimals:9}
Bilangan desimal dengan satu angka desimal mana saja yang terletak
tepat di antara $3$ dan $4$? Sebutkan semuanya. Ada berapa banyak?
\end{exercise}

\begin{exercise}[$\star\star$]\label{exo:g4:decimals:10}
Tono berkata: ``$0.25$ lebih besar daripada $0.5$ karena $25$ lebih
besar daripada $5$''. Tunjukkan kesalahannya dengan persegi seratus
(berapa petak yang diwarnai masing-masing?).
\end{exercise}

\begin{exercise}[$\star\star$]\label{exo:g4:decimals:11}
Tiga lompatan seorang pelari terukur $3.6$ m, $3.58$ m dan $3.65$ m.
Urutkan lompatannya dari yang terpendek ke yang terpanjang. Dua yang
mana yang berbeda tepat $\frac{5}{100}$ meter?

\end{exercise}
